\chapter{Results, Validation, Limitations and Future Work}
\label{chap:validation}

\section{Consolidated result}

The project now forms one inspectable chain: supplied CSV, causal preprocessing, frozen model artifacts, Python inference, versioned PostgreSQL records and a two-page DirectQuery dashboard. The strongest evidence is not simply that each component runs. It is that the model advantages and failures, source boundaries and interface responsibilities can all be traced to a specific artifact.

\begin{table}[htbp]
\centering\scriptsize
\caption{Consolidated quantitative evidence}
\label{tab:results-summary}
\begin{tabularx}{\textwidth}{@{}p{0.27\textwidth}p{0.26\textwidth}Y@{}}
\toprule
Evidence & Result & Correct interpretation \\
\midrule
Ridge TEST & MAE 0.001069; RMSE 0.001403; $R^2$ 0.824530; $n=165$ & Untouched synthetic chronological tail; percentage-point errors. \\
Persistence baseline & MAE 0.001459; RMSE 0.001848 & Ridge improves MAE 26.8\% and RMSE 24.1\%. \\
Residual audit & Bias 0.000066; Durbin--Watson 1.827; lag-1 autocorrelation 0.084 & Limited short-lag structure on this TEST sequence; not a plant stability claim. \\
One-Class SVM TEST & 2.654\% total flags; 60.25\% labelled-row response; 2.258\% normal-period flags & Process-novelty response on synthetic labelled episodes. \\
Episode response & Steam 82.0\%; airflow 79.3\%; feed surge 0\% & Sensitivity varies by departure type; the loading episode needs further work. \\
Full analytical replay & 237,600 rows in 2.037 s; about 116,626 rows/s & Batch model/feature computation only; database and Power BI excluded. \\
Database smoke test & 1,500 rows; 9/47 ms average/maximum cycle latency & Bounded integration evidence, not production load qualification. \\
Dashboard demo & 34,560 rows; 8 injections; 7/8 episode response & Jury presentation coverage only. \\
\bottomrule
\end{tabularx}
\source{Frozen artifacts and reproducible report audit; detailed evidence in Chapters~\ref{chap:models}--\ref{chap:powerbi}.}
\end{table}

\section{Verification performed on the final repository}

The current Python suite passes 77 tests plus 8 subtests. It covers timestamp and cadence contracts, causal alignment, feature order, artifact inference, replay behaviour, diagnosis rules, persistence, database helpers and presentation utilities. The Power BI report validator resolves 114 field references across 86 visuals against 303 model fields in six tables, with a 1600$\times$900 canvas.

A live PostgreSQL contract check also passes for all six DirectQuery views. At the time of the final check, the database contained a short 41-row current replay fragment, one laboratory row and no anomaly/contributor event rows. Empty optional event views are valid in a normal interval. This live check proves column compatibility and connectivity at that moment; it is not a substitute for the larger analytical or smoke-test results.

\begin{table}[htbp]
\centering\small
\caption{Validation layers and their remaining boundary}
\label{tab:validation-layers}
\begin{tabularx}{\textwidth}{@{}p{0.27\textwidth}p{0.32\textwidth}Y@{}}
\toprule
Layer & What passed & What it does not prove \\
\midrule
Unit and contract tests & 77 tests + 8 subtests & Plant representativeness or operator acceptance. \\
Report semantic validation & 86 visuals and all 114 field references resolve & Query latency under deployment concurrency. \\
Live SQL / Tabular Model Definition Language (TMDL) check & All six views match the semantic model & Long-duration resilience or gateway behaviour. \\
Model holdout & Chronological soft-sensor and synthetic-scenario anomaly audit & Accuracy on new OCP grades, seasons or real failures. \\
Presentation replay & Full-screen two-page flow with deterministic events & Independent model validation. \\
\bottomrule
\end{tabularx}
\end{table}

\section{What I consider demonstrated}

The work demonstrates that sparse laboratory values can be preserved while a faster process stream supports advisory inference. It shows that residence-aware alignment and strictly previous laboratory context can be implemented without obvious future leakage. It also shows a measurable soft-sensor improvement over a realistic persistence baseline.

The anomaly result is deliberately narrower. The frozen model responds strongly to two held-out disturbance families and does not flag the feed surge. What is demonstrated is an auditable novelty-to-guidance workflow, not comprehensive fault classification.

Finally, the complete software chain is coherent. The exact notebook artifacts reach the runtime, event records are idempotent and versioned, SQL views match the semantic model, and the dashboard preserves measurement/prediction semantics. This is an Industrial Engineering contribution because data quality, timing, interpretation and operator authority are treated as one system.

\section{Material limitations}

\begin{itemize}
\item \textbf{Synthetic source.} The supplied file supports reproducible prototype testing, but it cannot establish performance across real grades, seasons or operating modes.
\item \textbf{No live plant interface.} PCS 7 screens informed the analysis; CSV replay substitutes for historian and laboratory acquisition.
\item \textbf{Small supervised evidence.} The dense stream has 1,589,760 rows, but only 1,103 usable moisture targets. Label cadence, not process-row count, governs the supervised evidence.
\item \textbf{Scenario-dependent novelty detection.} The detector reacts to steam and airflow departures but needs representative plant history and confirmed events for calibration.
\item \textbf{Demonstration injections.} The two-day dashboard fork improves presentation coverage but cannot add credibility to model validation; it is disclosed and kept separate.
\item \textbf{Incomplete physical balance.} Wet-cake moisture, dry-solids flow, gas humidity/enthalpy and heat loss are unavailable, so thermal features remain proxies.
\item \textbf{Plant tag confirmation.} The canonical vacuum field follows the supplied convention; its polarity and pressure reference require confirmation.
\item \textbf{Prototype timing.} Five-second event/replay cadence and local performance do not qualify historian, network, gateway, concurrent Power BI or failover performance.
\item \textbf{No alarm or control authority.} Statistical thresholds are not approved operating, alarm, trip or safety limits, and no actuator command exists.
\end{itemize}

\section{Stage-gated industrialization roadmap}

\begin{figure}[htbp]
\centering
\begin{tikzpicture}[node distance=5mm]
\node[flow,minimum width=120mm,text width=113mm] (prototype)
  {\textbf{CURRENT PROTOTYPE}\\synthetic replay $\rightarrow$ Python inference $\rightarrow$ PostgreSQL views $\rightarrow$ Power BI};
\node[data,minimum width=120mm,text width=113mm,below=7mm of prototype] (gate1)
  {\textbf{GATE 1 -- data and semantics}\\approved historian/lab extract; tags, units, clocks, campaigns, shutdowns and quality flags reconciled};
\node[data,minimum width=120mm,text width=113mm,below=of gate1] (gate2)
  {\textbf{GATE 2 -- shadow validation}\\chronological retraining; baseline comparison; confirmed-event review; drift and nuisance-event burden};
\node[data,minimum width=120mm,text width=113mm,below=of gate2] (gate3)
  {\textbf{GATE 3 -- operational qualification}\\cybersecurity, gateway/load, recovery, alarm terminology, usability, ownership and rollback};
\node[flow,draw=OCPGreen,minimum width=120mm,text width=113mm,below=of gate3] (future)
  {\textbf{POSSIBLE ADVISORY SERVICE}\\read-only supervision with accepted accuracy, timing, stale-state and human-factor criteria};
\draw[dashedarr] (prototype)--(gate1); \draw[dashedarr] (gate1)--(gate2); \draw[dashedarr] (gate2)--(gate3); \draw[dashedarr] (gate3)--(future);
\end{tikzpicture}
\caption{Evidence gates from the current prototype to a possible industrial advisory service.}
\source{Author's work.}
\label{fig:futurearchitecture}
\end{figure}

\FloatBarrier

The next step should be data, not control. An approved read-only extract should cover normal campaigns, grade changes, shutdowns, maintenance and instrument-quality states, with laboratory sample and result times distinguished. The soft sensor should be retrained and compared by campaign against persistence and simple process baselines. The prototype currently provides point predictions; formal uncertainty estimation could be evaluated later using representative plant data.

For anomaly detection, confirmed events and normal operating modes are more valuable than another isolated hyperparameter search. A shadow review can label whether an event was actionable, expected, caused by data quality or unresolved. Those labels would support threshold selection, alternative semi-supervised models and a credible nuisance-event measure.

Only after data, model and service qualification should OCP consider a persistent advisory deployment. Any later recommendation layer would still require explicit operator authority and approved procedures. Autonomous control is outside the present roadmap.

\section{Personal contribution summary}

My main contribution was to connect disciplines that are often presented separately. I converted process observations into a defensible data contract, implemented causal alignment, compared the soft sensor against a practical baseline, exposed the anomaly model's missed scenario, built the runtime and database boundary, and shaped the dashboard around operator verification. I also added the evidence audit and report traceability needed to defend both the result and its limits. The remaining work is not hidden technical polish; it is the industrial evidence required before the prototype can leave a controlled environment.
